
\section{Comparar familias sin convertirlas en recetas}
En aluminio, las series 2xxx y 7xxx suelen relacionarse con endurecimiento por precipitación y elevadas resistencias en estados concretos; las 6xxx ofrecen otras combinaciones de fabricación y propiedades. La familia no determina una resistencia única ni garantiza soldabilidad, resistencia a corrosión o tenacidad. Un estado de tratamiento que maximiza una propiedad puede penalizar otra. Las aleaciones Al--Li añaden posibilidades de menor densidad y modificación del módulo, junto a exigencias de procesado y comportamiento que deben evaluarse en el producto.

En titanio, las proporciones y morfologías de fases alfa y beta dependen de composición e historia térmica. Las propiedades no se deducen únicamente del nombre de la fase: tamaño, distribución y textura importan. Se valoran resistencia específica y compatibilidad térmica, pero también fabricación, comportamiento de superficies y posibles daños por contacto. No se interpreta «resistente a corrosión» como inmunidad en todo ambiente.
\section{Discos y álabes: diferentes compromisos}
Un disco debe soportar esfuerzos asociados a rotación, gradientes térmicos, transitorios y concentraciones, con una población de defectos y una vida controladas. Un álabe caliente añade exposición severa a temperatura, ambiente y vibración y puede usar refrigeración y recubrimientos. Una microestructura adecuada para uno no es automáticamente la mejor para el otro.

Las superaleaciones de níquel combinan matriz y fases de refuerzo estables en un intervalo de servicio; composición, fracción, morfología y estabilidad de precipitados influyen en fluencia y fatiga. El procesado de disco o álabe busca microestructuras diferentes. Las superaleaciones de cobalto también ofrecen combinaciones útiles en aplicaciones calientes, pero no se sustituyen por las de níquel usando solo la temperatura de fusión. El comportamiento de oxidación, corrosión, fabricación y unión completa la comparación.
\section{Aceros, inoxidables y coste de ciclo de vida}
Aceros de alta resistencia pueden resultar apropiados para ejes, engranajes o elementos muy cargados, según diseño. Esa resistencia puede acompañarse de mayor sensibilidad a entallas, ambiente o tratamientos. Los inoxidables reúnen familias microestructurales distintas; la pasividad depende del medio y de la condición superficial. No existe un inoxidable universalmente inmune a picaduras o corrosión bajo tensión.

El coste de adquisición es solo una parte: herramientas, inspección, reparación, disponibilidad de producto, tiempos de suministro y retirada de servicio afectan al coste de ciclo de vida. Una preselección técnicamente prometedora puede descartarse si no existe una ruta de fabricación y control reproducible para la geometría requerida.
\repaso{¿La designación 8090 identifica una superaleación de níquel?}{No: corresponde a una aleación de aluminio--litio. La identificación de familia se verifica antes de discutir propiedades; un número de designación aislado no basta para asignar estado o admisibles.}
